\newcommand{\projecttitle}{Projet 3 - Parallélisation d’un~algorithme de distance d’édition}
\newcommand{\projectauthors}[1]{Romain Despoullains#1Nicolas Marano#1Sylvain Rowe#1Cloe Vautard}
\hypersetup{pdftitle={\projecttitle},pdfauthor={\projectauthors{, }},
  pdfsubject={CHPS1001 - Éléments de bioinformatique : utilisation du HPC dans la biologie}}

\newcommand{\doctitle}[1]{%
  \thispagestyle{empty}
  \begingroup
  \centering\setlength{\parskip}{0pt}
  \vspace*{8mm}
  % Les marges propres aux fichiers officiels sont retirées avant l’alignement.
  \includegraphics[trim=28.8bp 27.8bp 28.8bp 28.4bp,clip,height=17mm]{assets/logos/urca.png}%
  \hspace{11mm}%
  \includegraphics[trim=26bp 1.45bp 26bp 1bp,clip,height=16mm]{assets/logos/chps.jpg}\par
  \vspace{7mm}
  {\large Université de Reims Champagne-Ardenne\par}
  \vspace{2mm}
  Master 2 Calcul Haute Performance et Simulation (CHPS)\par
  Année universitaire 2026/2027\par
  \vfill
  \begin{minipage}{.90\linewidth}
    \centering
    \rule{\linewidth}{.4pt}\par\vspace{5mm}
    {\fontsize{25}{30}\selectfont\bfseries\projecttitle\par}
    \vspace{3mm}
    {\large\itshape Parallel Edit Distance for Genomic Sequences\par}
    \vspace{4mm}\rule{\linewidth}{.4pt}\par
  \end{minipage}\par
  \vspace{11mm}
  {\large #1\par}
  \vspace{10mm}
  {\bfseries Étudiants\par}
  \vspace{2mm}
  {\large\projectauthors{\\}\par}
  \vfill
  {\bfseries UE CHPS1001\par}
  Éléments de bioinformatique : utilisation du HPC dans la biologie\par
  \vspace{4mm}
  25 novembre 2026\par
  \vspace*{7mm}
  \endgroup
  \clearpage
}
